\documentclass[10pt]{article}
\usepackage[margin=1in]{geometry}
\usepackage{graphicx,booktabs,amsmath,hyperref,xcolor,listings}
\usepackage{caption}
\hypersetup{colorlinks=true,linkcolor=blue,urlcolor=blue}
\lstset{basicstyle=\ttfamily\footnotesize,breaklines=true,frame=single,
        backgroundcolor=\color{gray!8}}
\title{\textbf{Recreating ``I Know What You See'': A Power Side-Channel Attack on a
BNN Accelerator, Ported to CW305 + ChipWhisperer-Lite}\\
\large Reproduction, Hardware Realisation, and an Analysis of Why Measured-Silicon
Recovery Fails}
\author{Reproduction study}
\date{June 2026}

\begin{document}
\maketitle

\begin{abstract}
We reproduce the attack of Wei \emph{et al.} (ACSAC'18), which recovers the input image
of an FPGA binarized-CNN (BNN) accelerator from the power consumption of its first
convolution layer. We retarget the design from the original Spartan-6 / SAKURA-G /
2.5\,GHz-oscilloscope bench to a CW305 Artix-7 board captured with a ChipWhisperer-Lite.
The full pipeline is realised: a binarized MNIST CNN is trained, the first-layer
line-buffer convolution is implemented in HLS and verified \emph{bit-exact on silicon},
and the §6 (background detection) and §7 (power-template) attacks are implemented. In
\emph{simulation} the attack reproduces the paper (3$\times$3 template recognition
$0.935$ vs.\ paper $0.898$). On \emph{measured} silicon, recovery fails: the per-cycle
power correlates with the image only at $\rho\!\approx\!0.1$. We show through an
exhaustive sweep---four convolution micro-architectures (including a $16\times$ leakage
amplifier), two sampling regimes, and feature/gain/averaging variations---that the
limitation is the \emph{measurement bench}, not the implementation. We explain precisely
why the simulation succeeds while the hardware does not, compare the two benches, and
survey alternative routes (deep-learning SCA, EM, external high-rate capture, design
changes) that could recover a similar result.
\end{abstract}

\section{Introduction}
The target paper~\cite{wei2018} attacks the \emph{first} convolution layer of an
FPGA BNN accelerator. Because that layer's line buffer streams the raw image one pixel per
clock cycle, its dynamic power leaks pixel values; the authors recover MNIST digits with up
to $89.8\%$ recognition. The accelerator is the open \texttt{bnn-fpga} design of Zhao
\emph{et al.}~\cite{zhao2017}; the BNN methodology is BinaryNet~\cite{hubara2016}.

Our goal: reproduce this on a \emph{ChipWhisperer-Lite (CW1173)} capturing a
\emph{CW305 Artix-7 (XC7A100T)} target, using Vivado / Vivado~HLS~2016.4. This bench
differs fundamentally from the paper's (Sec.~\ref{sec:compare}), which turns out to be
decisive.

\section{Background}
\textbf{Line buffer (attack surface).} For a $K\!\times\!K$ kernel the conv unit holds
$K\!-\!1$ line buffers plus a $K\!\times\!K$ window of registers; each cycle one pixel
shifts in and one output is produced (paper Fig.~9c; \texttt{bnn-fpga} Sec.~III). The
moving pixels make the per-cycle switching data-dependent.

\textbf{Two attacks.} \emph{Background detection} (§6, passive): cycles whose window is
uniform (background) switch little $\Rightarrow$ low power; thresholding the per-cycle
power yields a silhouette. \emph{Power template} (§7, active): profile per-cycle power for
several kernels, then match an unknown image's per-cycle ``power feature vector'' against
the template to recover pixel values.

\section{Recreated setup}\label{sec:setup}
\textbf{Training (\texttt{training/}).} A 4-layer binarized CNN on MNIST, first layer
$64$ kernels, $K\!\in\!\{3,5\}$, implemented with a self-contained BinaryNet STE
(\texttt{bnn\_layers.py}; \texttt{larq} is incompatible with Keras~3). Test accuracy
$93\%$/$98.3\%$ ($3\!\times\!3$/$5\!\times\!5$); a BinaryNet MLP golden classifier reaches
$98.8\%$ for the recognition metric. The $64$ binarized kernels are exported
(\texttt{layer1\_kernels.npy}).

\textbf{FPGA conv (\texttt{hls/bnn\_conv1.cpp}, \texttt{cw305/}).} \texttt{bnn-fpga} is
SDSoC-only and cannot build with our toolchain, so the first-layer convolution is
re-implemented in HLS as a streaming line buffer (one pixel/cycle, II$=1$), wrapped in the
stock CW305 USB-register shell (\texttt{cw305\_usb\_reg\_fe.v}, \texttt{clocks.v},
\texttt{cw305.xdc}) by \texttt{cw305\_bnn\_top.v} / \texttt{cw305\_reg\_bnn.v}. The host
writes the image and one kernel, pulses \texttt{GO}, and the design asserts the
ChipWhisperer trigger for the whole convolution. The hardware output is \emph{bit-exact}
to a NumPy golden for every image/kernel tested (\texttt{host/cw305\_func\_check.py}).

\textbf{Capture.} The CW305 PLL clocks the conv; its clock is fed to the CW-Lite so the
ADC samples \emph{synchronously} (Fig.~\ref{fig:p0}). The amplified shunt (X4) feeds the
CW-Lite \textsc{measure} input.

\section{Attacks (\texttt{attack/})}
The §5 power-extraction front-end is replaced by a synchronous per-cycle reduction
(\texttt{power\_extract.py}); §6 uses an Otsu threshold (\texttt{attack\_background.py});
§7 builds a KD-tree template and reconstructs with a faithful Algorithm~2
(\texttt{power\_template.py}). Identical code runs on \emph{simulated} traces
(\texttt{trace\_sim.py}, a Hamming-distance CMOS power model) and on \emph{measured}
traces (\texttt{run\_on\_hardware.py}); only the trace source differs.

\section{Result 1 -- Simulation reproduces the paper}
With a Hamming-distance power model of the line buffer, the attack recovers the inputs at
paper level: $3\times3$ template recognition $\mathbf{0.935}$ (paper $0.898$), $5\times5$
$0.785$ (paper $0.790$); background recognition $0.77$/$0.655$ (paper $0.816$/$0.646$).
Fig.~\ref{fig:sim} shows recovered digits.

\begin{figure}[t]\centering
\includegraphics[width=\linewidth]{figs/sim_recovery.png}
\caption{Simulation: input recovered from the \emph{modelled} power. The template attack
restores recognisable digits (recognition $0.935$, paper $0.898$).}
\label{fig:sim}
\end{figure}

\section{Result 2 -- The bench works, but measured recovery fails}
The capture path is validated end-to-end with the stock AES bitstream
(Fig.~\ref{fig:p0}). The BNN conv produces a clean, triggered, synchronous trace
(Fig.~\ref{fig:real}). Yet running §6/§7 on \emph{measured} traces recovers nothing
(template recognition $\approx 0.15$ vs.\ $0.90$): the per-cycle power barely tracks the
image (Fig.~\ref{fig:corr}), $\rho\!\approx\!0.1$, far below the $\sim$0.5 needed.

\begin{figure}[t]\centering
\includegraphics[width=\linewidth]{figs/p0_aes.png}
\caption{P0: stock AES power trace on the CW305 -- the CW-Lite capture path (USB,
synchronous clock lock, trigger, amplified shunt) is validated.}
\label{fig:p0}
\end{figure}

\begin{figure}[t]\centering
\includegraphics[width=\linewidth]{figs/real_trace.png}
\caption{Measured layer-1 BNN convolution power trace (one kernel, $\mathrm{avg}=100$).
Triggered, synchronous, $4$ samples/cycle over $900$ cycles -- but the per-cycle dynamic
power is dominated by data-\emph{independent} activity.}
\label{fig:real}
\end{figure}

\begin{figure}[t]\centering
\includegraphics[width=\linewidth]{figs/corr_fail.png}
\caption{Measured per-cycle power does \emph{not} follow the true pixel content
(top), and the scatter (bottom) shows $\rho\!\approx\!0.1$ -- the attack cannot separate
foreground from background or match templates.}
\label{fig:corr}
\end{figure}

\section{Why simulation works but silicon does not}\label{sec:why}
The simulation uses an \emph{idealised per-cycle power model}: $p[c]$ is exactly the
Hamming distance of the line-buffer datapath between cycles $c\!-\!1$ and $c$, with additive
white noise. By construction $p[c]$ is a clean, perfectly-aligned function of the window
pixels, so §6/§7 recover the image.

Real silicon violates every idealisation:
\begin{enumerate}
\item \textbf{Per-cycle power is not observable at 4 samples/cycle.} AC-coupling and the
power-distribution-network capacitance \emph{smear} each cycle's transient into its
neighbours (the very effect the paper's §5 removes). The paper sampled at \textbf{2.5\,GHz}
($\sim$thousands of samples/cycle) and \emph{curve-fit the RC transient of each cycle} to
recover clean per-cycle power. With CW-Lite's \textbf{4 samples/cycle} that extraction is
impossible.
\item \textbf{The data-dependent fraction is tiny.} On the large XC7A100T, a sub-mW conv's
pixel-dependent switching is buried under clock-tree, static and I/O power. The paper used
a small Spartan-6 on SAKURA-G, an SCA-isolated board with a dedicated amplified power tap.
\end{enumerate}
The simulation models leakage \emph{as the paper's bench could measure it}; CW-Lite cannot
measure it that way.

\section{The fix attempts -- and the decisive negative result}
We aligned the implementation ever closer to the paper and tried to force the leakage:
\begin{itemize}
\item conv micro-architecture: (v1) preload + muxed reads, (v2) addressed line buffer,
(v3) true shift-register line buffer (paper Fig.~9c), (v4) a \textbf{$16\times$ parallel
leakage amplifier} (distinct unmergeable lanes all carrying the pixel stream);
\item sampling: synchronous $4\times$ and \textbf{asynchronous $21\times$/cycle}
oversampling (paper §5 style);
\item features (magnitude vs.\ Hamming distance), alignment scans, gain $30$--$50$\,dB,
averaging $50$--$500$.
\end{itemize}
All four micro-architectures are bit-exact on silicon. \textbf{Every variant produced the
same $\rho\!\approx\!0.1$--$0.2$} (Fig.~\ref{fig:wall}). Crucially, multiplying the
data-dependent switching $16\times$ \emph{did not change} the correlation: the bottleneck is
per-cycle \emph{extraction}, not leakage magnitude.

\begin{figure}[t]\centering
\includegraphics[width=0.85\linewidth]{figs/corr_wall.png}
\caption{Every design lever -- including a $16\times$ leakage amplifier -- lands at the
same $\sim$0.1 correlation floor. The wall is the measurement bench.}
\label{fig:wall}
\end{figure}

\section{Result 3 -- Deep-learning SCA, the strongest practical route, also fails}
\label{sec:dl}
Because a neural network can combine many weakly-leaking features non-linearly, we
implemented a profiled DL-SCA (\texttt{attack/dl\_sca.py}): for each pixel we gather the
trace samples of the $K\!\times\!K$ convolution cycles whose window contains it ($9$ kernels
$\times 4$ samples $\times 9$ cycles $=324$ features) and train an MLP to predict the pixel,
profiling on known images and recovering held-out ones. Per-pixel foreground/background
separability creeps up with data ($0.67\!\rightarrow\!0.71\!\rightarrow\!0.74$ for
$90\!\rightarrow\!300\!\rightarrow\!1000$ profiling images) but \emph{recognition never
leaves chance} ($0.067,\,0.14,\,0.095$; i.e.\ $\approx\!0.10$), and the recovered images
remain noise (Fig.~\ref{fig:dl}). The $0.74$ per-pixel ceiling---$26\%$ of pixels
wrong---is far below what a recognisable digit needs. The strongest practical attack, given
$1000$ profiling traces, still does not recover the input on this bench, which is the
cleanest evidence that the limit is the measurement, not the attack.

\begin{figure}[t]\centering
\includegraphics[width=\linewidth]{figs/dl_recovery.png}
\caption{Profiled deep-learning SCA on measured CW305 traces ($300$ profiling images).
Per-pixel accuracy $0.71$ and recognition $0.14$ reflect weak statistical bias only -- the
recovered images carry no digit structure.}
\label{fig:dl}
\end{figure}

\section{Paper bench vs.\ ours}\label{sec:compare}
\begin{table}[h]\centering\small
\begin{tabular}{lll}
\toprule
 & \textbf{Paper (Wei \emph{et al.})} & \textbf{This work} \\
\midrule
Target FPGA   & Spartan-6 LX75 (small)   & Artix-7 XC7A100T (large) \\
Board         & SAKURA-G (SCA-isolated)  & CW305 \\
Power tap     & dedicated amplified      & X4 amplified shunt \\
Oscilloscope  & Tektronix 2.5\,GHz       & CW-Lite 105\,MS/s \\
Samples/cycle & $\sim$thousands (async)  & 4 (sync) \\
§5 extraction & RC-transient curve fit   & not possible (4 samp/cyc) \\
Accelerator   & \texttt{bnn-fpga} (SDSoC)& HLS line buffer (bit-exact) \\
Recovery      & up to 89.8\%             & $\approx$15\% (sim: 93.5\%) \\
\bottomrule
\end{tabular}
\caption{The decisive differences are the oscilloscope bandwidth (samples/cycle) and the
SCA-isolated small target.}
\end{table}

\section{Alternatives to recover a similar result}\label{sec:alt}
Ranked by how much they keep of the original experiment:
\begin{enumerate}
\item \textbf{Deep-learning / profiled SCA (same hardware) -- implemented, Sec.~\ref{sec:dl}.}
With up to $1000$ profiling traces, recognition stays at chance ($\approx\!0.10$) and
recovery is noise; per-pixel separability plateaus at $0.74$. On present evidence it is
capped by the per-cycle extraction limit rather than by data, so this same-bench route does
not reach the paper.
\item \textbf{TVLA / Welch $t$-test first.} A fixed-vs-random pixel $t$-test
quantifies \emph{whether and where} per-cycle leakage exists, setting an upper bound on any
attack and validating capture before investing in (1).
\item \textbf{Better capture, same target.} An external GHz oscilloscope on the CW305 X4
SMA (or a ChipWhisperer-Husky, larger buffer + higher rate) restores the paper's
samples/cycle and enables §5 extraction; a CW501 differential probe and an external
low-noise VCC-INT supply raise SNR.
\item \textbf{EM instead of power.} An H-field probe over the convolution region localises
the signal and often beats global power SNR for a small core.
\item \textbf{Shrink the system-to-conv power ratio.} Deploy on a small FPGA (e.g.\
CW305-A35 or a Spartan-class board), clock-gate everything but the conv, and clock very low
so an external scope captures the transient -- approaching the paper's conditions.
\item \textbf{Change the leakage model deliberately.} Widen the pixel datapath / add a
register stage whose Hamming weight encodes the pixel onto a wide bus (a design change
within the line-buffer spirit), so the leak is large and easy to model.
\item \textbf{Reframe the deliverable.} The attack \emph{algorithm} is validated against
the paper in simulation; on this bench the realistic on-silicon demonstration is a simpler,
higher-leakage target (e.g.\ the AES key recovery the CW305 is designed for), with the BNN
input-recovery shown in simulation.
\end{enumerate}

\section{Conclusion}
We faithfully reproduced the attack's \emph{algorithm} (paper-level recovery in simulation)
and built the \emph{full hardware pipeline} (a bit-exact line-buffer conv on the CW305 with
working synchronous capture). Measured-silicon recovery does not reach the paper because the
CW305+CW-Lite bench cannot extract clean per-cycle power for a sub-mW convolution: proven by
a $16\times$ leakage amplifier leaving the correlation unchanged. Faithful real-silicon
recovery needs the paper's class of bench (SCA-isolated small target + GHz scope). Even a
profiled deep-learning attack with $1000$ traces (Sec.~\ref{sec:dl}) recovers only noise
(recognition at chance), confirming the floor is fundamental to this bench rather than to
any one attack. The faithful path forward is better acquisition (GHz scope / Husky, EM
probe, SCA-isolated small target); the existing CW305+CW-Lite setup is sufficient to
\emph{validate the design and algorithm} but not to \emph{measure} the leakage the attack
needs. All code, traces and figures are in the repository.

\begin{thebibliography}{9}\small
\bibitem{wei2018} L.~Wei, B.~Luo, Y.~Li, Y.~Liu, Q.~Xu, ``I Know What You See: Power
Side-Channel Attack on Convolutional Neural Network Accelerators,'' ACSAC 2018.
\bibitem{zhao2017} R.~Zhao \emph{et al.}, ``Accelerating Binarized Convolutional Neural
Networks with Software-Programmable FPGAs,'' FPGA 2017.
\bibitem{hubara2016} I.~Hubara, M.~Courbariaux \emph{et al.}, ``Binarized Neural
Networks,'' NeurIPS 2016.
\bibitem{timon2018} B.~Timon, ``Non-Profiled Deep Learning-based Side-Channel Attacks,''
IACR ePrint 2018/196.
\bibitem{dlsca} Practical Deep-Learning-based Power Side-Channel Attacks, arXiv:1907.02674.
\end{thebibliography}
\end{document}
